\section{研究背景与问题提出}

半导体是数字经济、先进制造和国家经济安全的重要基础。近年来，围绕先进计算、半导体制造设备和关键集成电路的出口管制不断加强，全球半导体贸易网络面临更强的政策不确定性 \citep{us_bis_controls}。对中国而言，进口来源是否集中、关键来源是否发生替代、不同产品风险是否一致，直接关系到产业链供应链韧性评估。

已有讨论常将半导体供应链风险概括为单一方向的“脱钩”或“断链”，但从统计建模角度看，贸易结构变化可能同时受到出口管制、全球产业周期、中国需求扩张和第三方产能调整影响。因此，本文不把美国出口管制视为唯一解释变量，而是将其作为重要政策背景，重点研究中国半导体相关产品进口来源结构的演化事实和风险暴露。

本文围绕三个问题展开：第一，2018 年以来，中国半导体相关产品进口来源结构是否发生明显变化；第二，不同产品的来源替代路径和风险暴露是否存在差异；第三，能否构建一个可解释、可复现的进口供应链风险评价指标，为后续政策监测提供量化依据。

本文的贡献主要体现在三点。其一，从单一产品扩展到设备类和集成电路类四个 HS6 产品，避免用单一产品概括整个半导体贸易结构。其二，构建 SIRI 风险指数，把来源集中、政策暴露、替代不足和结构波动纳入统一框架。其三，在固定效应回归不显著的情况下，明确区分描述性结构变化和严格因果识别，保持结论边界清楚。
